\section{Limitations}
\label{sec:limits}

Three load-bearing limits sit at the primitive layer.

\paragraph{No live cross-process federation.} The multi-party setup runs in-process: two kernels share a clock and exchange canonical bytes through library calls rather than the network. Cross-process federation and partition reconciliation sit in the deployment trajectory the related polity-layer construction outlines (anonymized for review).

\paragraph{Single-vendor key custody.} Bilateral DSSE carries party-independence only when the two signing keys are held by independent operators. If both kernels run under one custody domain the construction reduces to single-key signing, and operational discipline rather than cryptography carries the independence claim. A kernel-independence attestation primitive is named as future work.

\paragraph{Observability gap.} The receipt log is the audit channel: every admission and denial appears as a canonical receipt with a rejection code. Live observability (denial-rate alerts, schema-mismatch fingerprinting, non-leaking telemetry) is not specified.

\paragraph{Intra-lease replay state.} The primitive checks lease expiry against verifier time and lease epoch against verifier revocation state. It does not, by itself, consume an envelope or remember that a still-fresh envelope has already been presented. Preventing repeated use inside the same lease window requires a verifier-side seen-set, a receipt-graph continuation rule that consumes each continuation once, or an equivalent deployment rule. Without that state, the primitive rules out stale lease replay but not same-lease replay.

\paragraph{Side-channel attacks on the verifier.} The strict verifier is treated as a constant-time idealization. In practice, timing leaks along the signature-verification path, cache-resident key material, and microarchitectural side effects during canonical-bytes comparison can leak information about the trust-store membership decision or the byte-level canonical encoding. The construction does not defend against side-channel attacks on the verifier's runtime; UncoreBleed~\cite{uncoreBleed2026} and contemporary side-channel literature against EdDSA verifiers and TEE-resident comparators indicate that a hardware-isolated deployment must independently mitigate against such leaks. The contribution here is the wire-format admission predicate, not a constant-time verifier, and the receipt-graph evidence captures the verifier's accept-reject verdict rather than the byte-level execution trace an adversary observes.

\paragraph{Malicious receiving kernel.} The construction defends the sender's claim that a particular envelope was admitted under particular bytes; it does not defend against a receiver that controls its own trust store, suppresses denial receipts, or admits envelopes outside the declared predicate set. A receiving kernel that injects attacker-controlled keys into its allowlist can admit envelopes its declared policy should refuse; one that drops valid envelopes silently denies service. Such behavior is a constitutional violation rather than a cryptographic break: auditors with access to the receipt log detect the discrepancy by replay against canonical bytes, and a verifier that refuses to produce its receipt log at all is outside the scope of in-band detection. External assurance that a verifier evaluates the predicate it claims to evaluate is a kernel-attestation problem, addressed as future work.

\paragraph{Cryptographic-suite migration.} Ed25519 signatures and SHA-256 subject digests are not post-quantum. A quantum-equipped adversary breaks Ed25519 signing-key recovery directly and weakens the SHA-256 collision-resistance margin; both impact the admission predicate at the binding-tuple and signature-slice layers. The related polity-layer construction develops a salted-hash plus PQC-signature migration path (anonymized for review), and the bilateral admission primitive inherits that path. Deployed verifiers will need re-anchoring against PQC-signed envelopes; the wire format does not currently version-tag for cipher-suite agility, and adding such a tag (so a verifier can distinguish a v1 Ed25519/SHA-256 envelope from a successor profile without ambiguity) is part of the cryptographic-rotation work the related polity-layer construction outlines (anonymized for review).

Polity-level questions (citizenship, amendment, Hart-condition-(a) framing, trajectory-invariant theorems, sensor-state attestation) sit one layer above this primitive and are addressed in the related polity-layer construction (anonymized for review).
